\section{Messages are probability densities}
\label{sec:grid}

Unlike discrete-state models, here each latent variable
$a_k\in\mathbb R$. Consequently, a belief-propagation message
is a non-negative function of a continuous variable, rather
than a finite-dimensional vector. The message updates therefore
involve integrals over $\mathbb R$, which cannot in general be
evaluated in closed form for a non-Gaussian transition kernel.

We approximate these integrals numerically by discretising the
state space. We introduce a uniform grid
\begin{equation}
    g_i=-A+ih,\qquad i=0,\ldots,N_g-1,
    \qquad h=\frac{2A}{N_g-1},
\end{equation}
on $[-A,A]$, together with trapezoidal quadrature weights
$w_i$. The transition kernel and observation likelihood are
then evaluated on the grid:
\begin{equation}
    K_{ij}=K(g_i\mid g_j),
    \qquad
    (\ell_k)_i=\ell_t(x_k\mid g_i).
\end{equation}
Each integral in the forward and backward recursions becomes
a weighted matrix--vector product. In particular, the
unnormalised updates are
\begin{align}
    f_{k+1}
    &\propto
    \ell_k\odot K(w\odot f_k),\\
    b_k
    &\propto
    K^\top(w\odot\ell_{k+1}\odot b_{k+1}),
\end{align}
where $\odot$ denotes elementwise multiplication and
$w=(w_0,\ldots,w_{N_g-1})^\top$.
The messages are normalised after each update to ensure
numerical stability. The posterior mean is approximated by
\begin{equation}
    \widehat m_k
    =
    \frac{\sum_i w_i g_i(f_k)_i(b_k)_i}
         {\sum_i w_i(f_k)_i(b_k)_i},
\end{equation}
and substituted into the Tweedie identity to obtain the score.

The resulting computational cost is
$O(LN_g^2)$, with $N_g$ grid points and sequence length $L$.
Although this introduces a discretization error, its
accuracy can be controlled by refining the grid and enlarging
the integration domain, and against the Gaussian reference.

\paragraph{Why not a Gaussian approximation?}
A natural alternative would be to approximate each message
by a Gaussian distribution, retaining only its mean and
variance. This is exact for the Gaussian chain: Gaussian
likelihoods and transition kernels are closed under the
multiplication and marginalisation operations of belief
propagation. It is not exact for the Laplace chain, however.
A Gaussian closure discards the non-Gaussian features of the
messages and effectively recovers the score of the
covariance-matched Gaussian model.

This distinction matters most at low noise, where the
Laplace-chain score exhibits a non-linear, saturating
dependence on the observations. As shown in
Fig.~\ref{fig:grid-vs-gaussian}, the Gaussian approximation
can therefore incur substantial score error, whereas grid
belief propagation closely tracks the reference solution.

This should also be distinguished from approximate message
passing (AMP) and Thouless--Anderson--Palmer (TAP) methods.
Those methods are motivated by dense, high-dimensional
systems, where many weak interactions justify approximate
Gaussian message updates, often with an Onsager correction.
Our factor graph is instead a finite-degree chain and a tree:
ordinary belief propagation is already exact at the
continuous level. We therefore discretise the messages
rather than introduce a Gaussian approximation.